\subsection{Any length metric is isospectral to a non-positively
  curved Riemannian metric}
\label{sec:non-pos}

Use Hamenstädt's result: Liouville currents of negatively curved
Riemannian metrics are dense in currents \cite[Lemma 2.8]{Ham99:Cocycles}

Bound diameter following MathOverflow answer
\url{https://mathoverflow.net/questions/346811/can-the-diameter-be-controled-by-the-injectivity-radius-and-the-volume}:
\begin{lemma}
  For any non-positively curved metric $g$ on a surface~$\Sigma$,
  \[
    \diam(g) \le \frac{8\vol(g)}{\systole(g)}.
  \]
\end{lemma}
(Similar results hold in any dimension, and without the assumption of
non-positive curvature, replacing systole with injectivity radius.)

Combine with Gromov compactness. Show that $\pi_1$ doesn't change in
limit, again using positive systole.
\DMGtodo[inline]{Gromov compactness requires uniform bound on curvature. but these metrics need not have uniform bounded sectional curvature, for example it could converge to -infinity as it converged to a tree?}
Unique if result is negatively curved, not in general (who did this?)
\DMGtodo[inline]{For negatively curved Riemannian metrics it was Hersonsky-Paulin in the paper whose argument we used to prove countable additivity in our argument. But rigidity is true in more generality than just negatively curved Riemannian: for non-positively curved Riemannian was proved by Constantine in 2018. For non-positively curved Euclidean cone metrics the latest result is by Viveka, Chris Leininger and Chandrika Sadanand. I would say the first page of this preprint contains the lastest account of the history of this problem \url{https://arxiv.org/pdf/2104.12176.pdf}. One thing to note is that if we can actually show that our metric in the isospectral result is Riemannian then because of rigidity there is a unique such metric (it's uniquely determined by the lengths). So interestingly we are picking up a distinguished non-positively curved Riemannian metric... Kind of a uniformization procedure.}

\begin{remark}
  Gromov-Hausdorff limit can be nasty in general, for instance if we
  approximate a closed curve or measured lamination. Systole is zero
  in these cases.
\end{remark}

\DPTtodo[inline]{BIPP 2019 prove that, for intersection with a current
$\mu$, systole positive $\Leftrightarrow$ $\mu$ is filling. Seems very
relevant}


\begin{proof}
Given any (length) metric $\rho$ on $S$, consider $\ell_{\rho}$ its function length of free homotopy classes of curves.
By our Theorem~\ref{thm:intersect}, there exist a geodesic current $\mu_{\rho}$ so that
\[
i(\mu_{\rho},C)=\ell_{\rho}(C)
\]
for any closed curve $C$.
By \cite[Lemma~2.18]{Ham99:Cocycles}, there is a sequence of Liouville currents $L_{m_i}$ associated to negatively curved metrics on $S$ so that $L_{m_i} \to \mu_{\rho}$.

\end{proof}